% Generated draft: requires mathematical and source review.
\documentclass[11pt]{article}
\usepackage[margin=25mm]{geometry}
\usepackage{fontspec}
\setmainfont{Latin Modern Roman}
\usepackage{amsmath,amssymb,mathtools,unicode-math}
\setmathfont{Latin Modern Math}
\usepackage{enumitem,needspace,xurl,hyperref}
\setlength{\parindent}{0pt}
\setlength{\parskip}{5pt}
\setlength{\emergencystretch}{3em}
\newcommand{\N}{\mathbb{N}}
\newcommand{\Z}{\mathbb{Z}}
\newcommand{\Q}{\mathbb{Q}}
\newcommand{\R}{\mathbb{R}}
\newcommand{\C}{\mathbb{C}}
\newcommand{\F}{\mathbb{F}}
\newcommand{\A}{\mathbb{A}}
\newcommand{\PP}{\mathbb{P}}
\newcommand{\E}{\mathbb{E}}
\newcommand{\eps}{\varepsilon}
\newcommand{\dd}{\,\mathrm{d}}
\begin{document}
\section*{UnsolvedMath ID: 8800016}
\textbf{Source problem code:} catalogue entry
\subsection{Definitions and mathematical statement}
In a finite-field setting \(q=p^a\), a deformation method for Frobenius computation uses a family of curves \(C_t\) connecting a target curve to a starting fiber \(C_{t_0}\) for which the Frobenius action on a chosen cohomology space is easy to compute. For a basis \(v_1,\ldots,v_m\), a Frobenius matrix records \(F(v_j)=\sum_i A_{ij}v_i\), with linear or semilinear conventions determined by the chosen Frobenius and coefficient field.

Kedlaya's task is to find \emph{useful deformations of nondegenerate curves together with easy starting matrices} for this computation. Both the deformation family and the computable initial Frobenius data are required; an abstract deformation without explicit starting data does not supply the requested construction.

One common meaning of nondegenerate in this context is Newton-polygon nondegeneracy for a Laurent equation \(f(x,y)=0\). If \(f_\tau\) is its face polynomial, the condition is that, for every face \(\tau\), the equations
\[
 f_\tau=x\frac{\partial f_\tau}{\partial x}
 =y\frac{\partial f_\tau}{\partial y}=0
\]
have no simultaneous solution in \((\overline{\F}_q^{\times})^2\). The source excerpt does not provide this definition, so it remains a marked conventional interpretation rather than an imposed replacement.

The desired curve classes, deformation base, allowed singular fibers, cohomology theory and basis, Frobenius power, \(p\)-adic precision, and efficiency requirements are all unspecified. ``Useful'' and ``easy'' do not define numerical complexity targets. The question asks for constructions, not a universal optimality theorem or a claim that one family handles every nondegenerate curve.
\subsection{Short English statement}
Find deformation families for nondegenerate curves that aid Frobenius computation, together with explicitly accessible Frobenius matrices at starting fibers. The intended nondegeneracy and computational criteria need the missing source context.
\subsection{Sources}
\begin{itemize}
\item Local UnsolvedMath registry, permanent ID 8800016, source code catalogue entry; attributed to Kiran Kedlaya; no proposal year supplied; updated 2026-08-06.
\item IPAM, \emph{Workshop on Number Theory and Cryptography --- Open Problems} (2006), attachment 6220, PDF pages 36--38, logical slide 19: \url{https://www.ipam.ucla.edu/programs/workshops/workshop-i-number-theory-and-cryptography-open-problems/?tab=schedule}. Registry access date 2026-07-29; PDF and legacy-PPT visual audit. Provenance reports 19 surviving attachments, 17 PDFs and two legacy PPT decks, with repetitions and solved definitions excluded. Present status and redistribution terms require review.
\item The supplied partial-progress summary mentions \(p\)-adic deformation methods, sparse hyperelliptic starting equations, Dwork-style techniques and inverse Teichm\"uller starting points, without exact references, precise algorithms, or a verified complete selection theory.
\end{itemize}
\end{document}
